\SourceLeaf{29}
\fancyhead[L]{\small GENERAL FUNDAMENTAL LAW OF MECHANICS.}
\pdfbookmark[0]{On a New General Fundamental Law of Mechanics}{mechanik}
\begin{center}
ON A NEW\par
\bigskip
GENERAL FUNDAMENTAL LAW OF MECHANICS.\par
\medskip
\rule{.16\textwidth}{.3pt}
\end{center}
\expandafter\def\csname HistoricalNoteText2\endcsname{Allow me to remark here, however, that I find unsatisfactory the way another great geometer tried to prove \textsc{Huyghens}’ law for the extraordinary refraction of light in doubly refracting crystals by means of the principle of least action. Indeed, whether this principle is applicable depends essentially on the principle of conservation of vis viva. According to the latter, the speeds of the moving material points are determined solely by their positions, with no possible influence from the direction of motion. Yet the attempt just mentioned assumes that the direction does have an influence. It seems to me that under the emission theory of light all efforts to connect the phenomena of double refraction with the general laws of dynamics are bound to fail as long as the particles of light are treated simply as points.}
\par
As is well known, the principle of virtual velocities turns all of statics into a mathematical problem, and \textsc{d’Alembert}’s principle for dynamics in turn reduces dynamics to statics. It follows from the nature of the subject that there \textit{can} be no new fundamental principle for the study of motion and equilibrium that is not already contained in substance in those two principles and derivable from them. Even so, it hardly follows that every new principle is worthless. It will always be interesting and instructive to find a new and advantageous way of looking at the laws of nature, whether it makes one particular problem or another easier to solve or brings out some special appropriateness. The great geometer who built the edifice of mechanics so brilliantly on the foundation of the principle of virtual velocities did not disdain to give \textsc{Maupertuis}’ principle of least action greater precision and generality. This is a principle that can sometimes be used to great advantage\HistoricalNote{2}.
